%% =====================================================================
%%  Unit 18 — What did you do last weekend?  (Raconter son week-end)
%%  Pistes : 18.01–18.04 dialogues · 18.05–18.06 vocabulaire · 18.07–18.08 grammaire
%%           18.09 prononciation · 18.10–18.13 exercices (textes en annexe)
%% =====================================================================

\newcommand{\otIllusUnXVIII}{\otIllusDuo{watched/104/86/15/0.10,brunch/190/92/20/0.09,stayed/122/14/13/0.08,
  did/150/8/11/0.08,visited/196/22/12/0.07,played/100/50/12/0.07}{emma}{What did you do?}{lea}{I tried brunch!}}

\unite[What did you do\\last weekend?]{18}{What did you do last weekend?}{Raconter son week-end}{otIllusUnXVIII}

\situation{Épisode 18}{Lundi matin, fin janvier. Au studio, Emma veut tout savoir sur le week-end de
Léa : elle a découvert le \en{brunch} du dimanche ! Un peu plus tard, Mr.\ Carter fait le point avec
elle sur un projet urgent.}

\begin{objectifs}
\item raconter ce que vous avez fait, au prétérit ;
\item former le prétérit des verbes réguliers et l'écrire correctement ;
\item poser des questions avec \en{Did…?} et \en{What did…?} et y répondre ;
\item dire ce que vous n'avez pas fait, avec \en{didn't} ;
\item prononcer la terminaison \en{-ed} : \phon{t}, \phon{d} ou \phon{ɪd}.
\end{objectifs}

\auprogramme
  {Raconter des événements passés}
  {Le prétérit des verbes réguliers\par\vspace{0.6mm}\en{did}, \en{didn't}}
  {Verbes réguliers, week-end : 26 expressions}
  {La terminaison \en{-ed} : \phon{t}, \phon{d}, \phon{ɪd}}
  {Le \en{brunch} du dimanche}

\semaine{casque/otNavy/Dialogues/Écouter et lire les deux scènes,
  bulle/otGreenDark/Vocabulaire/Apprendre les 26 expressions,
  blocs/otViolet/Grammaire/Le prétérit et did,
  sons/otOrangeDark/Sons et culture/La terminaison -ed et le brunch,
  crayon/otNavy/Exercices/12 exercices corrigés,
  coche/otGreenDark/Bilan/L'essentiel et le test express}

\clearpage
%% ---------------------------------------------------------------------
\rubrique{ecouter}{Raconter son week-end}

\begin{dialogue}{1}{What did you do last weekend?}{Lundi, 9 heures, au studio. Emma pose son manteau
  et rejoint Léa près de la machine à café.}{18.01}{18.02}
\dit{emma}{Morning, Léa! \hl{What did you do} last weekend?\note{1}}{Bonjour, Léa ! Qu'est-ce que tu as fait le week-end dernier ?}
\dit{lea}{On Saturday, I \hl{stayed} home. I \hl{cleaned} my apartment and I \hl{watched} a movie.\note{2}}{Samedi, je suis restée chez moi. J'ai fait le ménage et j'ai regardé un film.}
\dit{emma}{And on Sunday?}{Et dimanche ?}
\dit{lea}{On Sunday, I \hl{tried} brunch for the first time! Sam \hl{invited} me.\note{3}}{Dimanche, j'ai testé le brunch pour la première fois ! Sam m'a invitée.}
\dit{emma}{Oh, lovely! \hl{Did you like} it?\note{4}}{Oh, super ! Ça t'a plu ?}
\dit{lea}{\hl{Yes, I did}! I \hl{ordered} pancakes with maple syrup. Delicious!}{Oui ! J'ai commandé des pancakes au sirop d'érable. Délicieux !}
\dit{emma}{\hl{Did} Sam \hl{order} pancakes too?}{Sam a pris des pancakes, lui aussi ?}
\dit{lea}{\hl{No, he didn't}. He \hl{wanted} eggs and bacon.\note{5}}{Non. Il voulait des œufs au bacon.}
\end{dialogue}

\begin{comprendre}
\pointnote{1} \en{\textbf{What did you do?}} Au prétérit, les questions se forment avec \en{did} +
  sujet + base verbale : « qu'est-ce que tu as fait ? ».
\pointnote{2} \en{\textbf{I stayed, I cleaned, I watched}} Le prétérit des verbes réguliers : base
  verbale + \en{-ed}. C'est la même forme pour toutes les personnes. Il traduit le passé composé
  (« j'ai regardé ») et souvent l'imparfait.
\pointnote{3} \en{\textbf{I tried}} \en{try} → \en{tried} : après une consonne, le \en{y} devient
  \en{ied}. Dans \en{invited}, la terminaison se prononce \phon{ɪd}.
\pointnote{4} \en{\textbf{Did you like it? — Yes, I did.}} La réponse courte reprend \en{did}.
\pointnote{5} \en{\textbf{No, he didn't.}} \en{didn't} (= \en{did not}) sert à la négation : \en{He
  didn't order pancakes}. Après \en{did} et \en{didn't}, le verbe reste à la base verbale.
\end{comprendre}

\begin{dialogue}{2}{Did you call Mr. Davis?}{Lundi, 11 heures, dans le bureau de Mr.\ Carter. Le projet
  pour Lumen Books doit être prêt cette semaine.}{18.03}{18.04}
\dit{carter}{Léa, the Lumen Books project. \hl{Did you call} Mr.\ Davis?}{Léa, le projet Lumen Books. Vous avez appelé Mr.\ Davis ?}
\dit{lea}{Yes, I did. I \hl{called} him at nine.}{Oui. Je l'ai appelé à neuf heures.}
\dit{carter}{Great. \hl{Did} he \hl{like} the new logo?}{Très bien. Le nouveau logo lui a plu ?}
\dit{lea}{Yes, he \hl{loved} it! But he \hl{asked} for a different color.}{Oui, il l'a adoré ! Mais il a demandé une autre couleur.}
\dit{carter}{OK. \hl{Did you change} it?}{D'accord. Vous l'avez modifié ?}
\dit{lea}{Yes. I \hl{finished} it \hl{ten minutes ago}, and I \hl{emailed} it to him.\note{6}}{Oui. Je l'ai terminé il y a dix minutes, et je le lui ai envoyé par e-mail.}
\dit{carter}{Perfect! You \hl{worked} very hard last week. Thank you, Léa.}{Parfait ! Vous avez beaucoup travaillé la semaine dernière. Merci, Léa.}
\dit{lea}{You're welcome. Oh, and Mr.\ Davis \hl{invited} us to lunch on Friday!}{Je vous en prie. Ah, et Mr.\ Davis nous a invités à déjeuner vendredi !}
\end{dialogue}

\begin{comprendre}
\pointnote{6} \en{\textbf{ten minutes ago}} Avec un moment passé et terminé (\en{yesterday}, \en{last
  week}, \en{ago}, \en{at nine}), l'anglais emploie toujours le prétérit : \en{I finished it ten minutes
  ago}.
\end{comprendre}

%% ---------------------------------------------------------------------
\rubrique{mots}{Les mots de l'unité}

Voici les 26 expressions de l'unité. Pour chaque verbe, apprenez la base verbale et le prétérit
ensemble, à voix haute : votre oreille retiendra le son de la terminaison.

\begin{vocabulaire}{Verbes réguliers}{18.05}
\categorie{Terminaison \phon{t}}
\mot[watch, watched]{watch · watched}{wɑːtʃ · wɑːtʃt}{regarder}
\mot[walk, walked]{walk · walked}{wɑːk · wɑːkt}{marcher, se promener}
\mot[cook... cooked]{cook · cooked}{kʊk · kʊkt}{cuisiner}
\mot[dance, danced]{dance · danced}{dæns · dænst}{danser}
\mot[stop, stopped]{stop · stopped}{stɑːp · stɑːpt}{(s')arrêter}
\categorie{Terminaison \phon{d}}
\mot[stay, stayed]{stay · stayed}{steɪ · steɪd}{rester}
\mot[clean, cleaned]{clean · cleaned}{kliːn · kliːnd}{nettoyer}
\mot[call, called]{call · called}{kɑːl · kɑːld}{appeler}
\mot[love, loved]{love · loved}{lʌv · lʌvd}{adorer}
\mot[try, tried]{try · tried}{traɪ · traɪd}{essayer, goûter}
\mot[study, studied]{study · studied}{ˈstʌdi · ˈstʌdid}{étudier}
\categorie{Terminaison \phon{ɪd}}
\mot[visit, visited]{visit · visited}{ˈvɪzɪt · ˈvɪzɪtɪd}{visiter}
\mot[want, wanted]{want · wanted}{wɑːnt · ˈwɑːntɪd}{vouloir}
\mot[start, started]{start · started}{stɑːrt · ˈstɑːrtɪd}{commencer}
\end{vocabulaire}

\begin{vocabulaire}{Le week-end}{18.06}
\categorie{Raconter}
\mot{last weekend}{ˌlæst ˈwiːkend}{le week-end dernier}
\mot{on Saturday morning}{ɑːn ˌsætɚdeɪ ˈmɔːrnɪŋ}{samedi matin}
\mot{What did you do?}{ˌwʌt dɪd juː ˈduː}{Qu'est-ce que tu as fait ?}
\mot{Did you have a nice weekend?}{dɪd juː ˌhæv ə ˌnaɪs ˈwiːkend}{Tu as passé un bon week-end ?}
\mot{I didn't do much.}{aɪ ˌdɪdənt duː ˈmʌtʃ}{Je n'ai pas fait grand-chose.}
\mot{It was fun!}{ɪt wəz ˈfʌn}{C'était sympa !}
\categorie{Activités}
\mot{stay home}{ˌsteɪ ˈhoʊm}{rester à la maison}
\mot{relax}{rɪˈlæks}{se détendre}
\mot{a series}{ə ˈsɪriːz}{une série (télé)}
\mot{brunch}{brʌntʃ}{le brunch}
\mot{pancakes}{ˈpænkeɪks}{des pancakes}
\mot{maple syrup}{ˌmeɪpəl ˈsɪrəp}{le sirop d'érable}
\end{vocabulaire}

%% ---------------------------------------------------------------------
\rubrique{grammaire}{Le prétérit des verbes réguliers}

\otIntertitre[otViolet]{1 · La forme : base verbale + \en{-ed}}
\begin{center}\otTableau
\begin{tabular}{@{}l@{\hspace{8mm}}>{\enOn}l<{\enOff}@{\hspace{8mm}}l@{}}
\otTH{Forme} & \otTH{Exemple} & \otTH{Français}\\\midrule
affirmation      & I / you / he / she / we / they watch\textbf{ed} TV.       & j'ai regardé…\\
négation         & She \textbf{didn't} watch TV.                              & elle n'a pas regardé…\\
question         & \textbf{Did} you watch TV? · What \textbf{did} you watch?  & tu as regardé… ?\\
réponses courtes & Yes, I \textbf{did}. · No, he \textbf{didn't}.             & oui / non\\\bottomrule
\end{tabular}\end{center}

\begin{regle}[Avec \en{did}, une seule marque du passé]
Comme \en{do} au présent, \en{did} porte la marque du passé : le verbe qui suit reste à la base
verbale. On dit \en{Did you like it?} et \en{He didn't call}, jamais \en{Did you liked} ni \en{He
didn't called}. Bonne nouvelle : \en{did} est le même pour toutes les personnes.
\end{regle}

\blocaudio{18.07}{Écoutez et répétez}{Six phrases au prétérit.}

\otIntertitre[otViolet]{2 · L'orthographe du \en{-ed}}
\begin{center}\otTableau
\begin{tabular}{@{}l@{\hspace{8mm}}>{\enOn}l<{\enOff}@{}}
\otTH{Règle} & \otTH{Exemples}\\\midrule
cas général : on ajoute \en{-ed}                 & work → worked · play → played · stay → stayed\\
\en{-e} final : on ajoute seulement \en{-d}       & love → loved · dance → danced · invite → invited\\
consonne + \en{y} : \en{-ied}                     & try → tried · study → studied\\
voyelle + consonne finale : on la double          & stop → stopped · plan → planned\\\bottomrule
\end{tabular}\end{center}

\blocaudio{18.08}{Écoutez et répétez}{Une question avec \en{Did…?}, puis sa réponse.}

\begin{attention}[Les pièges des francophones]
\fauxjuste{I watch TV yesterday.}{I watched TV yesterday.}
\fauxjuste{Did you liked it?}{Did you like it?}
\fauxjuste{She didn't called me.}{She didn't call me.}
\fauxjuste{I studyed English.}{I studied English.}
\fauxjuste{What you did last weekend?}{What did you do last weekend?}
\end{attention}

%% ---------------------------------------------------------------------
\rubrique[La terminaison -ed]{sons}{La terminaison \en{-ed} : \phon{t}, \phon{d} ou \phon{ɪd}}

Erreur classique : prononcer \en{-ed} comme « ède ». En réalité, on n'ajoute une syllabe que dans un cas
sur trois. Comme pour le \en{-s} du pluriel (unité 5), c'est le son final du verbe qui décide.

\begin{tcbraster}[raster columns=3,raster equal height=rows,raster column skip=3mm,raster before skip=2mm,raster after skip=3mm]
\otCartePrg{sons}{otOrangeDark}{\phon{t}}{après \phon{p}, \phon{k}, \phon{s}, \phon{ʃ}, \phon{tʃ} (sons « sans voix ») : \en{stopped}, \en{walked}, \en{danced}, \en{watched}}
\otCartePrg{sons}{otOrangeDark}{\phon{d}}{après une voyelle ou un son « avec voix » : \en{played}, \en{cleaned}, \en{called}, \en{loved}}
\otCartePrg{sons}{otOrangeDark}{\phon{ɪd}}{après \phon{t} ou \phon{d} seulement : une syllabe en plus ! \en{visited}, \en{wanted}, \en{needed}}
\end{tcbraster}

\begin{paires}{18.09}{Écoutez la base verbale, puis le prétérit, et répétez. Seuls \en{visited} et
\en{wanted} ont une syllabe de plus.}
\sonpaire{walk}{wɑːk}{\phon{t}}{walk\otH{ed}}{wɑːkt}{marché}
\sonpaire{watch}{wɑːtʃ}{\phon{t}}{watch\otH{ed}}{wɑːtʃt}{regardé}
\sonpaire{play}{pleɪ}{\phon{d}}{play\otH{ed}}{pleɪd}{joué}
\sonpaire{clean}{kliːn}{\phon{d}}{clean\otH{ed}}{kliːnd}{nettoyé}
\sonpaire{visit}{ˈvɪzɪt}{\phon{ɪd}}{visit\otH{ed}}{ˈvɪzɪtɪd}{visité}
\sonpaire{want}{wɑːnt}{\phon{ɪd}}{want\otH{ed}}{ˈwɑːntɪd}{voulu}
\end{paires}

%% ---------------------------------------------------------------------
\rubrique{culture}{Le brunch du dimanche}

\begin{culture}[\en{Breakfast} + \en{lunch} = \en{brunch}]
Le dimanche, entre 11 heures et 14 heures, les New-Yorkais vont « bruncher » entre amis. Au menu :
\en{pancakes} ou \en{waffles} (gaufres) au sirop d'érable, \en{eggs Benedict} (œufs pochés, sauce
hollandaise), \en{bagels}, bacon croustillant… et café à volonté (\en{free refills}). Dans les quartiers à
la mode, on fait souvent la queue (\en{a line}) une demi-heure !
\end{culture}

\begin{paire}
\begin{culture}[Le sirop d'érable]
Le \en{maple syrup} vient de la sève de l'érable à sucre, récoltée à la fin de l'hiver. Le Québec en
produit environ les trois quarts de la production mondiale ; aux États-Unis, c'est la spécialité du
Vermont.
\end{culture}
\begin{culture}[Le week-end américain]
Le samedi, on fait ses courses et ses \en{errands} (petites démarches), on accompagne les enfants au
match de \en{soccer}. Le dimanche, d'automne en hiver, beaucoup d'Américains regardent le football
américain à la télévision.
\end{culture}
\end{paire}

%% ---------------------------------------------------------------------
\rubrique{pratique}{Exercices}

\exercice[ecrire,niveau=1]{Au prétérit}{Écrivez le prétérit de chaque verbe.}
\begin{anglais}
\begin{questions}[4]
\item play → \trou[18mm]{played}
\item stop → \trou[18mm]{stopped}
\item try → \trou[18mm]{tried}
\item love → \trou[18mm]{loved}
\item study → \trou[18mm]{studied}
\item walk → \trou[18mm]{walked}
\item visit → \trou[18mm]{visited}
\item plan → \trou[18mm]{planned}
\end{questions}
\end{anglais}

\exercice[lire,ecrire,niveau=1]{Le week-end de Jack}{Lisez l'agenda de Jack, puis complétez les phrases
au prétérit.}
\begin{center}
\begin{tcolorbox}[enhanced,width=118mm,colback=white,colframe=otRule,boxrule=0.6pt,arc=2mm,
  left=4mm,right=4mm,top=3mm,bottom=3mm,
  title={\otSemi\fontsize{9}{11}\selectfont\enb{Jack's weekend}},coltitle=white,colbacktitle=otNavy,
  fonttitle=\otSemi,attach boxed title to top left={xshift=4mm,yshift=-2.5mm},boxed title style={arc=1mm,boxrule=0pt}]
\vspace{2mm}\renewcommand{\arraystretch}{1.4}
\begin{tabular}{@{}>{\otSemi\fontsize{9}{11}\selectfont\color{otNavy}\enOn}l<{\enOff}@{\hspace{6mm}}>{\enOn}l<{\enOff}@{\hspace{4mm}}c@{}}
Sat., 9 a.m. & clean the apartment & \otCheck\\
Sat., 2 p.m. & play hockey with Kate & \otCheck\\
Sat., 8 p.m. & watch a movie & \otCheck\\
Sun., 11 a.m. & brunch with Sam and Léa & \otCroix\\
Sun., 3 p.m. & visit Mom and Dad & \otCheck\\
Sun., 7 p.m. & cook dinner & \otCheck\\
\end{tabular}
\end{tcolorbox}
\end{center}
\begin{anglais}
\begin{questions}
\item On Saturday morning, Jack \trou[26mm]{cleaned} the apartment.
\item In the afternoon, he \trou[26mm]{played} hockey with Kate.
\item In the evening, he \trou[26mm]{watched} a movie.
\item On Sunday morning, he \trou[26mm]{didn't} go to the brunch.
\item In the afternoon, he \trou[26mm]{visited} his parents.
\item Then he \trou[26mm]{cooked} dinner.
\end{questions}
\end{anglais}

\exercice[ecrire,niveau=2]{À la forme négative}{Récrivez chaque phrase à la forme négative.}
\begin{questions}
\item I watched TV. \hfill \enb{\trou[66mm]{I didn't watch TV.}}
\item She called me. \hfill \enb{\trou[66mm]{She didn't call me.}}
\item We played tennis. \hfill \enb{\trou[66mm]{We didn't play tennis.}}
\item They visited Boston. \hfill \enb{\trou[66mm]{They didn't visit Boston.}}
\item He studied last night. \hfill \enb{\trou[66mm]{He didn't study last night.}}
\end{questions}

\exercice[ecrire,niveau=2]{Posez la question}{Écrivez la question au prétérit.}
\begin{questions}
\item \en{you} / \en{call Emma} \hfill \enb{\trou[66mm]{Did you call Emma?}}
\item \en{she} / \en{like the movie} \hfill \enb{\trou[66mm]{Did she like the movie?}}
\item \en{what} / \en{you} / \en{cook} \hfill \enb{\trou[66mm]{What did you cook?}}
\item \en{they} / \en{visit Boston} \hfill \enb{\trou[66mm]{Did they visit Boston?}}
\item \en{where} / \en{he} / \en{stay} \hfill \enb{\trou[66mm]{Where did he stay?}}
\end{questions}

\exercice[lire,niveau=2]{Le blog de Léa}{Lisez le message de Léa, puis indiquez si chaque phrase est
vraie (V) ou fausse (F).}
\begin{anglais}
\begin{tcolorbox}[otBase,colback=white,colframe=otRule,boxrule=0.5pt,fontupper=\fontsize{10.8}{14.4}\selectfont]
{\otSemi\fontsize{11}{13}\selectfont My first brunch!}\par\vspace{1.4mm}
Last Sunday, Sam invited me to brunch in Williamsburg. We arrived at eleven, but there was a line!
We waited for thirty minutes. I ordered pancakes with maple syrup and a big coffee. Sam wanted eggs
and bacon. We talked for two hours. I loved it! Brunch is my new favorite thing in New York.
\par\vspace{1mm}{\fontsize{8.6}{10}\selectfont\color{otMuted}\foreignlanguage{french}{\textit{Williamsburg} : un quartier de Brooklyn · \textit{a line} : une file d'attente · \textit{talk} : parler, discuter}}
\end{tcolorbox}
\begin{questions}[2]
\item Léa and Sam arrived at eleven. \hfill\trou[9mm]{V}
\item They didn't wait. \hfill\trou[9mm]{F}
\item Léa ordered pancakes. \hfill\trou[9mm]{V}
\item Sam wanted pancakes too. \hfill\trou[9mm]{F}
\item They talked for two hours. \hfill\trou[9mm]{V}
\item Léa didn't like brunch. \hfill\trou[9mm]{F}
\end{questions}
\end{anglais}

\exercice[ecrire,niveau=2]{Corrigez l'erreur}{Chaque phrase contient une erreur. Récrivez-la
correctement.}
\begin{questions}
\item I watch TV yesterday. \hfill \enb{\trou[66mm]{I watched TV yesterday.}}
\item Did you liked it? \hfill \enb{\trou[66mm]{Did you like it?}}
\item She didn't called me. \hfill \enb{\trou[66mm]{She didn't call me.}}
\item I studyed English. \hfill \enb{\trou[66mm]{I studied English.}}
\item We play tennis last week. \hfill \enb{\trou[66mm]{We played tennis last week.}}
\end{questions}

\exercice[ecouter,niveau=1,piste=18.10]{\phon{t}, \phon{d} ou \phon{ɪd} ?}{Écoutez chaque verbe au
prétérit, dit deux fois. Cochez le son de la terminaison.}
\begin{questions}[3]
\item \choix{\phon{t}}\quad\choix{\phon{d}}\quad\choix{\phon{ɪd}}
\item \choix{\phon{t}}\quad\choix{\phon{d}}\quad\choix{\phon{ɪd}}
\item \choix{\phon{t}}\quad\choix{\phon{d}}\quad\choix{\phon{ɪd}}
\item \choix{\phon{t}}\quad\choix{\phon{d}}\quad\choix{\phon{ɪd}}
\item \choix{\phon{t}}\quad\choix{\phon{d}}\quad\choix{\phon{ɪd}}
\item \choix{\phon{t}}\quad\choix{\phon{d}}\quad\choix{\phon{ɪd}}
\end{questions}

\exercice[ecouter,ecrire,niveau=2,piste=18.11]{Le week-end de Maya}{Sam demande à sa sœur comment
s'est passé son week-end. Écoutez et complétez.}
\begin{anglais}
\begin{questions}
\item On Saturday, Maya visited a \trou[30mm]{museum}.
\item Then she walked in \trou[30mm]{Central Park}.
\item On Sunday, she \trou[30mm]{stayed} home.
\item She \trou[30mm]{studied} for an exam.
\item In the evening, she watched a \trou[30mm]{series}.
\end{questions}
\end{anglais}

\exercice[ecouter,ecrire,niveau=2,piste=18.12]{Dictée à trous}{Emma et Jack parlent du week-end.
Écoutez et complétez.}
\begin{anglais}
\noindent\begin{tabular}{@{}l@{\hspace{6mm}}l@{}}
\otNomDit{emma} & Did you have a nice weekend, Jack?\\[1.8mm]
\otNomDit{jack} & Yes, I \trou{did}. I \trou{played} hockey with Kate.\\[1.8mm]
\otNomDit{emma} & Did you \trou{watch} the game on Sunday?\\[1.8mm]
\otNomDit{jack} & No, I \trou{didn't}. I \trou{visited} my parents.\\[1.8mm]
\otNomDit{emma} & Lovely! What did you \trou{do} there?\\[1.8mm]
\otNomDit{jack} & I \trou{cooked} dinner for them. They \trou{loved} it!\\
\end{tabular}
\end{anglais}

\exercice[ecrire,niveau=2]{Traduisez}{Traduisez en anglais.}
\begin{questions}[1][french]
\item Qu'est-ce que tu as fait le week-end dernier ?\lignes{1}
\item J'ai regardé un film et j'ai cuisiné.\lignes{1}
\item Elle ne m'a pas appelé.\lignes{1}
\item Tu as aimé le brunch ? — Oui. \textit{(réponse courte)}\lignes{1}
\item Nous avons visité Boston l'année dernière.\lignes{1}
\item Il a étudié hier soir.\lignes{1}
\item Ils ont dansé jusqu'à minuit.\lignes{1}
\item J'ai terminé il y a dix minutes.\lignes{1}
\end{questions}

\exercice[ecrire,niveau=3]{Mon week-end}{En six phrases, racontez votre dernier week-end : ce que vous
avez fait samedi et dimanche, avec qui, et ce que vous n'avez pas fait.}
\lignes{6}

\exercice[ecouter,parler,niveau=3,piste=18.13]{À vous de parler !}{Léa vous pose des questions sur votre
week-end. Répondez \textbf{à voix haute} pendant chaque pause.}
\begin{center}\otTableau
\begin{tabular}{@{}>{\otSemi\fontsize{8.6}{10}\selectfont\color{otNavy}}l@{\hspace{3mm}}>{\enOn}p{70mm}<{\enOff}@{\hspace{6mm}}p{78mm}@{}}
 & \otTH{Léa dit\ldots} & \otTH{Vous\ldots}\\\midrule
1. & Did you have a nice weekend? & répondez par une réponse courte.\\
2. & What did you do on Saturday? & racontez votre samedi.\\
3. & Did you stay home on Sunday? & répondez et précisez.\\
4. & Did you watch a movie? & répondez et dites lequel.\\
5. & Did you cook something nice? & répondez par une phrase complète.\\\bottomrule
\end{tabular}\end{center}

%% ---------------------------------------------------------------------
\rubrique{bilan}{Faites le point}

\begin{essentiel}
\begin{minipage}[t]{0.54\linewidth}
\otEssTitre{Les phrases clés}
\renewcommand{\arraystretch}{1.32}
\begin{tabular}{@{}>{\enOn\otSemi\fontsize{9.4}{11}\selectfont}l<{\enOff}@{\hspace{3mm}}>{\fontsize{9.2}{11}\selectfont\color{white!78!otNavy}}l@{}}
What did you do?           & Qu'est-ce que tu as fait ?\\
I watched a movie.         & J'ai regardé un film.\\
I didn't call him.         & Je ne l'ai pas appelé.\\
Did you like it?           & Ça t'a plu ?\\
Yes, I did. / No, I didn't. & Oui. / Non.\\
It was fun!                & C'était sympa !\\
\end{tabular}
\end{minipage}\hfill
\begin{minipage}[t]{0.41\linewidth}
\otEssTitre{Le prétérit régulier}
{\fontsize{9.2}{11.6}\selectfont base + \en{-ed}, pour toutes les personnes · \en{tried}, \en{stopped},
\en{loved}\par}
\otEssTitre{\en{did} · \en{didn't}}
{\fontsize{9.2}{11.6}\selectfont \en{Did you…?} · \en{I didn't…} + base verbale, sans \en{-ed}\par}
\otEssTitre{Prononciation}
{\fontsize{9.2}{11.6}\selectfont \en{walked} \phon{t} · \en{played} \phon{d} · \en{visited}
\phon{ɪd}\par}
\end{minipage}
\end{essentiel}

\begin{jesais}
\item raconter mon week-end au prétérit ;
\item écrire correctement le prétérit des verbes réguliers ;
\item poser une question avec \en{Did…?} ou \en{What did…?} et y répondre ;
\item dire ce que je n'ai pas fait, avec \en{didn't} ;
\item prononcer \en{-ed} \phon{t}, \phon{d} ou \phon{ɪd} ;
\item présenter le \en{brunch} américain.
\end{jesais}

\clearpage
\begin{testexpress}
\qcm[b]{« J'ai regardé la télé » se dit :}{\en{I watch TV.}}{\en{I watched TV.}}{\en{I did watch TV.}}
\qcm[c]{Le prétérit de \en{study} est :}{\en{studyed}}{\en{studed}}{\en{studied}}
\qcm[a]{« Elle ne m'a pas appelé » se dit :}{\en{She didn't call me.}}{\en{She didn't called me.}}{\en{She not called me.}}
\qcm[b]{\en{Did you like it?} — Réponse affirmative :}{\en{Yes, I liked.}}{\en{Yes, I did.}}{\en{Yes, I do.}}
\qcm[c]{« Qu'est-ce que tu as fait ? » se dit :}{\en{What you did?}}{\en{What did you?}}{\en{What did you do?}}
\qcm[a]{Le prétérit de \en{stop} est :}{\en{stopped}}{\en{stoped}}{\en{stopt}}
\qcm[b]{Dans \en{visited}, \en{-ed} se prononce :}{\phon{t}}{\phon{ɪd}}{\phon{d}}
\qcm[c]{Dans \en{walked}, \en{-ed} se prononce :}{\phon{ɪd}}{\phon{d}}{\phon{t}}
\qcm[a]{« J'ai terminé il y a dix minutes » se dit :}{\en{I finished ten minutes ago.}}{\en{I finish ten minutes ago.}}{\en{I finished there is ten minutes.}}
\qcm[b]{Le \en{brunch}, c'est :}{un dîner du samedi}{un repas entre petit-déjeuner et déjeuner}{un goûter}
\end{testexpress}

\noindent{\fontsize{9.6}{12}\selectfont\color{otMuted}Niveau visé (CECRL, niveau A1) : « peut raconter
très simplement des événements passés ». Les corrigés de l'unité sont page~\pageref{cor-u18-1} ; les
textes des pistes d'exercices, page~\pageref{trans-18.10}.\par}
